\section{An explicit cutoff at size 100,000}
\label{sec:degree-block-cutoff}

\textbf{Theorem.} For every \(n\ge100{,}000\) and \(3\le k\le n\), the
\(n-2\) zeros of \(P_n=E_{n,k}/t\) are simple and negative. They strictly
interlace those of \(P_{n+1}\), with the next row supplying both outer
zeros. Zero is a simple common root of the unreduced rows.

Write \(S=\sin\theta\), \(C=\cos\theta\), \(t=-e^{-\pi\cot\theta}\) and
\(z_a=(C+iaS)^{-1}\) for odd \(a\), where \(0<\theta\le\pi/2\).
Set \(F_j=(\pi/S)^j E_{j,k}(t)/(j!(1-t)^j)\), \(B_j=\sum_a z_a^j\)
and \(m=n+1-k\). Let \(D(u)=\exp(\sum_{j\ge k}u^j/j)=\sum_r d_r u^r\).
At any old zero the insertion identity gives
\((n+1)F_{n+1}=SF'_n+B_kF_m\), with the prime denoting differentiation
in \(\theta\). Define
\[
 V_\pm=k\left(F_n,\frac{F'_n}{n}\pm\frac{B_kF_m}{nS}\right),
 \qquad U=2kd_n\bigl(\cos(n\theta),-\sin(n\theta)\bigr).
\]
If \(\|V_\pm-U\|_2<2kd_n\), then \(V_\pm\cdot U>0\) by Cauchy--Schwarz.
Thus \(F_n\) has sign \((-1)^j\) at \(\theta=j\pi/n\). At an old root
between phases \(j\pi\) and \((j+1)\pi\), both second coordinates have
sign \((-1)^{j+1}\). Their average gives \(F'_n\ne0\), and the plus
coordinate gives sign \((-1)^j\) for \(P_{n+1}\) because \(t<0\).

To prove the disk inequalities, put \(b_s=[u^s](1-u)(D(u)-1)\). The signed
pole expansion for \(M=n,m\) distinguishes one occupied pole \(a\)
and writes \(F_M=\sum_a[u^M](D(z_a u)-1)Q_a(u)\), where
\[
 Q_a(u)=\prod_{q\ne a}\left(1+\sum_{s\ge k}
     \frac{b_s z_q^{s-1}u^s}{i(q-a)S}\right).
\]
The constant term of \(Q_a\) gives \(d_M B_M\). In the old row, a mixed
term \(z_a^r h\), where \(T\) is the degree from the other poles, \(r=n-T\), and
\(h\) denotes the scalar from the other poles, enters the value and
normalized derivative through a real two-by-two triangular matrix.
Its diagonal entries are \(1\) and \((r/n)a|z_a|^2\). The absolute
off-diagonal entry is at most \((r/n)(a^2-1)S|z_a|^2\).
Writing \(\rho_a=|z_a|\), the latter cost times \(\rho_a^r\) is at
most \(|a|\rho_a^{r/2}/\sqrt n\) and vanishes at the principal poles.
The separate derivative of \(h\) is bounded by \(|h'|\rho_a^r/n\).
This estimates the two coordinates jointly.

Assign bin \(i=\lceil 4s/k\rceil\) to each primitive factor of degree
\(s\). A nonempty product in bin \(b\) has total degree \(T\) with
\(\max(1,b/5)\le T/k\le b/4\), since each factor has degree at least
\(k\). The bounds on bin mass follow from the atom \(b_k=1/k\), the
telescoping differences \(1/(s(s-1))\) for \(k<s<2k\), the jump of
size at most \(1/(4k^2)\) at \(2k\), and the block bound
\(|b_s|\le 1/(2j!k^2)\) for \(jk<s\le(j+1)k\), \(j\ge2\).
For a slab \(\ell\le n/k\le X\), truncate the bin polynomials at
\(5X\). Sum absolute primitive terms and their normalized angular
derivatives into nonnegative majorants \(L(v)\) and \(J(v)\), using
\(kS\ge A\), \(nS^2\ge 6/5\), and the positive kernels summed over signed odd indices
\((1+c(p^2-1))^{-96}\) given by Bernoulli's inequality. Allowing
repeated pole choices gives \(\exp L-1\) for the mixed terms and
\(J\exp L\) for their differentiated factors. In each bin the
remaining distinguished degree is bounded using the interval for
\(T\). The cycle probability bounds \(kd_r\le4/7\) for \(r\ge2k\)
and \(kd_n\ge547/1000\) for \(n\ge3k\) sharpen the comparison.
The two rational replays enclose the pole sums and compare the
resulting joint error, including the insertion term, against
\(2kd_n\). The least certified margin exceeds \(65716/10^6\).
\href{https://github.com/ProgramComputer/open-math-research/tree/main/topics/long-cycle-excedances/degree-blocks}{Detailed bounds and exact replays}

\begin{center}
\begin{tabular}{lll}
Minimum \(n\) & Slabs for \(x=n/k\) & Sufficient \(A\)\\ \hline
\(100{,}000\) & \([4,6], [6,10], [10,16], [16,25]\) & \(10\)\\
\(100{,}000\) & \([25,34], [34,42], [42,50]\) & \(9,8,27/4\) respectively\\
\(100{,}000\) & \([50,62], [62,75], [75,86]\) & \(7,7,27/4\) respectively\\
\(120{,}000\) & \([86,96]\) & \(7\)
\end{tabular}
\end{center}

The exterior theorem supplies the first \(\lfloor7\sqrt n/20\rfloor\)
ranks at each end. Its last root precedes the first middle sample,
where \(nS^2\ge6/5\). This inequality gives \(kS\ge
\sqrt{(6/5)n}/X\) on each slab, enough for the table through \(x=50\).
For the later slabs it gives \(kS\ge3\). When \(kS<A\) there,
\(kS^2<A^2/k\le1/25\), so the differentiated common-residue comparison
gives the same integer-phase signs and the same actual next-row
signs at old roots. The comparisons may therefore be used at
different points of one phase interval. Reciprocity and the degree
give exactly one old root in each remaining gap, with the central
root at \(-1\) for odd \(n\). The next-row signs and the endpoint
coefficients give strict interlacing.

The three-cycle theorem covers \(x<4\). For \(100{,}000\le n<120{,}000\)
the table reaches \(x=86\), and for \(120{,}000\le n<4{,}062{,}720\)
it reaches \(x=96\). At and above these respective joins
\(0.89x>47+2.5\log n\). The lower-height common-circle comparison
applies under this condition for \(k\ge400\), while its small-threshold
part applies to \(4\le k<400\) from size \(55{,}000\). The fixed
threshold theorem covers \(k=3\), and the preceding uniform cutoff
covers \(n\ge4{,}062{,}720\).
